\section{Lubin--Tate branches, jet recovery and Fleck congruences}\label{sec:lub}

Berger and Sprang recover bivariate Taylor coefficients from substitutions along finitely
many Lubin--Tate branches in their study of bounded functions on the
Schneider--Teitelbaum character variety: Lemma~4.2 of \cite{BergerSprang2025Integer} gives a
finite truncation by a compactness argument, and Proposition~4.4 there uses it to extract
coefficients. This section computes the least possible truncation. The answer is one plus
the pole order of an inverse interpolation denominator. Over the residue field it is a sum
of contact orders; modulo $\pi^r$ it grows linearly in $r$ with a slope read from the
deepest contact, up to a bounded, eventually periodic defect that vanishes exactly under a
base-$p$ carry condition. The optimal choice of directions is the same at every precision.
Specialized to the multiplicative group, the pole computation becomes a simultaneous
congruence for weighted Fleck sums. The last subsection treats a related invariant of
\cite[Definition~5.1]{BergerSprang2025Integer}, the coefficient depth $s_0(n)$, through an
exact valuation formula for powers of the Lubin--Tate logarithm.

All results of this section are local and algebraic. They use no hypothesis on zeros and
do not depend on \cref{src:thm:dgrh}. The bounded-functions conjecture that motivates
\cite{BergerSprang2025Integer} has been reported for quadratic extensions of $\Qp$
\cite{BergerSprangSteingart2026Bounded} and announced for all finite extensions
\cite{Colovic2026Ring}; the statements below neither rely on those results nor bear on
them, with the single exception of the finiteness of $s_0(n)$ recalled in
Section~\ref{lub:ssec:depth}.

\subsection{Branches, jets and the interpolation threshold}\label{lub:ssec:threshold}
\begingroup\raggedright\sources{grh_transfers_cycle_signings_and_lubin_tate; grh_transfers_lubin_tate_and_graph_towers; grh_prime_patterns_and_branch_precision; grh_method_transfers; grh_method_transfers_orthogonal_jets_2026-11-26}\endgroup

Let $K/\Qp$ be a finite extension with ring of integers $\OO$, uniformizer $\pi$, valuation
$v$ normalized by $v(\pi)=1$, residue field $k$ of size $q$, and put $e=v(p)$. Fix a
Lubin--Tate formal $\OO$-module for $\pi$ in a coordinate $X$, and write $[b](X)$ for
multiplication by $b\in\OO$. For $r\ge1$ put $A_r=\OO/\pi^r$, so that $A_1=k$. Let
$B\subset\OO$ be a finite set of $d\ge1$ distinct \emph{directions}, and for $b\in B$ put
\begin{gather*}
 u_b(X)=[b](X)\bmod\pi^r\in XA_r[[X]],\\
 P_b(X,Y)=\prod_{c\in B\setminus\set{b}}\bigl(Y-u_c(X)\bigr),\qquad
 \Delta_b(X)=P_b\bigl(X,u_b(X)\bigr).
\end{gather*}
Let $H\in A_r[[X,Y]]$ be unknown. The \emph{observation} on the branch $b$ at precision
$n_b$ is the residue $H(X,u_b(X))\bmod X^{n_b}$; precision $n$ retains exactly the
coefficients of $X^0,\dots,X^{n-1}$. The desired output is the \emph{jet} of $H$: the list of its
coefficients of total degree $<d$.

\begin{lemma}[Contact of two branches]\label{lub:lem:contact}
Let $b\ne c$ in $\OO$ and $h=v(b-c)$. In $\OO[[X]]$,
\[
 [b](X)-[c](X)=[\pi^h](X)\cdot U(X),\qquad U\in\OO[[X]]^\times ,
\]
and the reduction of $[\pi^h](X)$ modulo $\pi$ has order exactly $q^h$ in $X$. Consequently
the reductions of the two branches over $k$ meet to order $q^{v(b-c)}$, and $u_b-u_c$ is a
unit of $A_r((X))$ for every $r$.
\end{lemma}
\begin{proof}
For the formal group law $\mathcal F$, the series $\mathcal F(S,T)-S$ equals $T$ times a
power series with constant term $1$. Substituting $S=[c](X)$ and $T=[b-c](X)$ shows that
$[b]-[c]$ is a unit multiple of $[b-c]$. If $b-c=\pi^hw$ with $w\in\OO^\times$, then
$[b-c]=[w]\circ[\pi^h]$, and $[w](Z)$ is $Z$ times a unit because its linear coefficient is
$w$. The reduction of $[\pi]$ has order $q$, and orders multiply under composition.
Finally a series $X^sV+\pi V'$ with $V\in\OO[[X]]^\times$ is invertible in $A_r((X))$:
factor off $X^sV$ and expand the remaining factor $1+\pi X^{-s}V^{-1}V'$ in a geometric
series, which is finite because $\pi$ is nilpotent in $A_r$.
\end{proof}

By Lemma~\ref{lub:lem:contact}, $\Delta_b$ is a unit of $A_r((X))$. Define the \emph{inverse pole
order}
\begin{equation}\label{lub:eq:pole}
 t_b(r)=\min\set{t\ge0:\ X^t/\Delta_b\in A_r[[X]]}.
\end{equation}
For $d=1$ the empty product is $1$ and $t_b(r)=0$.

\begin{proposition}[Exact componentwise recovery threshold; \stU]\label{lub:prop:threshold}
The observations $H(X,u_b(X))\bmod X^{n_b}$, $b\in B$, determine the jet of every
$H\in A_r[[X,Y]]$ if and only if
\[
 n_b\ge1+t_b(r)\qquad\text{for every }b\in B.
\]
In particular the least common precision is $1+\max_{b\in B}t_b(r)$.
\end{proposition}
\begin{proof}
Put $P=\prod_{b\in B}(Y-u_b(X))$. Division by this monic distinguished polynomial in $Y$
over the complete local ring $A_r[[X]]$ gives a remainder $R$ of $Y$-degree $<d$. Since
every $u_b$ has zero constant term, $P\in(X,Y)^d$, so $H$ and $R$ have the same jet. Over
$A_r((X))$ Lagrange interpolation is valid because every $u_b-u_c$ is a unit:
\begin{equation}\label{lub:eq:lagrange}
 R(X,Y)=\sum_{b\in B}\frac{H(X,u_b(X))}{\Delta_b(X)}\,P_b(X,Y).
\end{equation}
If all observations vanish and $n_b\ge1+t_b(r)$ for every $b$, each quotient in
\eqref{lub:eq:lagrange} belongs to $XA_r[[X]]$. Since $P_b\in(X,Y)^{d-1}$, the remainder
belongs to $(X,Y)^d$ and the jet of $H$ vanishes.

For necessity fix $b$ and set $g_b=X^{t_b(r)}/\Delta_b\in A_r[[X]]$. Its constant
coefficient is nonzero: otherwise $g_b$ would be divisible by $X$, contradicting the
minimality of $t_b(r)$ (for $t_b(r)=0$ the claim is immediate from $g_b\Delta_b=1$). The
series
\[
 H_b(X,Y)=g_b(X)\,P_b(X,Y)
\]
has the nonzero coefficient $g_b(0)$ at $Y^{d-1}$, vanishes on every other branch, and
evaluates to $X^{t_b(r)}$ on the branch $b$. If $n_b\le t_b(r)$, all its observations
vanish while its jet does not, so recovery fails, even with exact data on every other
branch. The argument needs only $g_b(0)\ne0$, not that $g_b(0)$ be a unit, and so remains
valid in the presence of nilpotents in $A_r$.
\end{proof}

The rest of the analysis concerns $t_b(r)$ for one fixed $b\in B$, with $d\ge2$. Write
\begin{equation}\label{lub:eq:contacts}
\begin{gathered}
 h_c=v(b-c)\quad(c\in B\setminus\set{b}),\qquad
 \omega=\omega_b=\sum_{c\ne b}q^{h_c},\\
 \nu_j=\#\set{c\ne b: h_c=j},\qquad
 m_j=\#\set{c\ne b: h_c\ge j},
\end{gathered}
\end{equation}
and call the multiset $\set{h_c}$ the \emph{contact multiset} of $b$. In the Lubin--Tate
coordinate with
\begin{equation}\label{lub:eq:coordinate}
 [\pi](X)=\pi X+X^q
\end{equation}
define the polynomials
\begin{equation}\label{lub:eq:iterates}
\begin{gathered}
 \phi_0=X,\qquad \phi_h=\pi\phi_{h-1}+\phi_{h-1}^{\,q}=[\pi^h](X)\quad(h\ge1),\\
 \Phi_j=\phi_j/\phi_{j-1}=\pi+\phi_{j-1}^{\,q-1},\qquad D_j=(q-1)q^{j-1}\quad(j\ge1).
\end{gathered}
\end{equation}
Each $\Phi_j$ is monic of degree $D_j$, has constant term $\pi$, and is congruent to
$X^{D_j}$ modulo $\pi$. For a monic polynomial $Q\in\OO[X]$ with
$Q\equiv X^{\deg Q}\pmod\pi$, write $\tau_r(Q)$ for the least $n\ge0$ such that $X^n=0$ in
$A_r[X]/(Q)$; thus $\tau_r(1)=0$.

\begin{lemma}[Normal form of the denominator]\label{lub:lem:normal}
\begin{enumerate}[label=\textup{(\roman*)}]
\item In any Lubin--Tate coordinate, $\Delta_b$ is the product of a unit of $A_r[[X]]$ and
the reduction of $\prod_{c\ne b}[\pi^{h_c}](X)$, and $t_b(r)$ is the least $t$ for which
$X^t/\prod_{c\ne b}[\pi^{h_c}](X)\in A_r[[X]]$. This number is the same for all Lubin--Tate
formal $\OO$-modules for $\pi$ and all their integral coordinates. For fixed $K$ and $\pi$
it therefore depends only on $r$ and on the contact multiset, not on the unit parts of the
differences $b-c$.
\item In the coordinate \eqref{lub:eq:coordinate},
\begin{equation}\label{lub:eq:denominator}
 P^{(b)}(X):=\prod_{c\ne b}\phi_{h_c}(X)=X^{d-1}\prod_{j\ge1}\Phi_j(X)^{m_j},
 \qquad \deg P^{(b)}=d-1+\sum_{j\ge1}D_jm_j=\omega .
\end{equation}
\item If $Q\in\OO[X]$ is monic with $Q\equiv X^{\deg Q}\pmod\pi$, then
$\min\set{t:\ X^t/Q\in A_r[[X]]}=\tau_r(Q)$. In particular $t_b(r)=\tau_r(P^{(b)})$.
If $Q$ divides a second such polynomial $Q'$, then $\tau_r(Q)\le\tau_r(Q')$.
\end{enumerate}
\end{lemma}
\begin{proof}
(i) The first assertion is Lemma~\ref{lub:lem:contact}. Multiplication by a unit power series
preserves the least exponent of every nonzero Laurent series over $A_r$: the coefficient at
that exponent is multiplied by the unit constant term and remains nonzero, also when it is
nilpotent. This gives the description of $t_b(r)$. By the theorem of Lubin and Tate
\cite{LubinTate1965Formal}, any two Lubin--Tate formal $\OO$-modules for $\pi$ are
isomorphic over $\OO$; more generally let $\theta(X)=uX+\cdots$, $u\in\OO^\times$, be an
integral change of coordinate conjugating one formal module into the other. Then
$\theta(S)-\theta(T)=(S-T)\Theta(S,T)$ with $\Theta(0,0)=u$, so every transformed branch
difference is the old difference, after the invertible substitution
$X\mapsto\theta^{-1}(X)$, times a unit. This substitution is an automorphism of $A_r((X))$
preserving every lattice $X^jA_r[[X]]$, hence the least exponent of a Laurent series. The
total-degree filtration of $A_r[[X,Y]]$ is preserved as well.

(ii) The identity $\phi_j=\phi_{j-1}\Phi_j$ gives $\phi_h=X\prod_{j=1}^h\Phi_j$, and
$q^h=1+\sum_{j=1}^hD_j$.

(iii) Since $Q$ is a unit of $A_r((X))$, the condition $X^t/Q\in A_r[[X]]$ says that
$X^t\in QA_r[[X]]$. Division by the monic polynomial $Q$ gives the same remainder in the
polynomial ring $A_r[X]$ and, by the uniqueness in Weierstrass division, in $A_r[[X]]$. So
the condition says that $X^t$ vanishes in $A_r[X]/(Q)$. If $Q\mid Q'$, then $X^n\in(Q')$
implies $X^n\in(Q)$.
\end{proof}

Part (iii) yields a terminating exact algorithm for $t_b(r)$: start with $1$, multiply by
$X$, and reduce modulo $P^{(b)}$ and $\pi^r$ until the remainder is zero.

\subsection{The sharp pole bound and the carry criterion}\label{lub:ssec:pole}
\begingroup\raggedright\sources{grh_transfers_cycle_signings_and_lubin_tate; grh_transfers_lubin_tate_and_graph_towers; grh_prime_patterns_and_branch_precision; grh_method_transfers; grh_method_transfers_orthogonal_jets_2026-11-26}\endgroup

\begin{lemma}[The leading Newton face]\label{lub:lem:face}
For $a\ge1$ write $\phi_a=\sum_jc_{a,j}X^j$. The polynomial $\phi_a$ is monic of degree
$q^a$ and reduces to $X^{q^a}$ modulo $\pi$. Every nonzero lower coefficient satisfies
\[
 q^a-j\le D_a\,v(c_{a,j}).
\]
Among the lower coefficients equality occurs only at $j=q^{a-1}$, where
$c_{a,j}\equiv\pi\pmod{\pi^2}$. Moreover every exponent occurring in $\phi_a$ is congruent
to $1$ modulo $q-1$.
\end{lemma}
\begin{proof}
The case $a=1$ is immediate from \eqref{lub:eq:coordinate}. For the inductive step note
that $D_a=qD_{a-1}$ and use the recurrence \eqref{lub:eq:iterates}. In a nonleading product
monomial from $\phi_{a-1}^{\,q}$, the total degree deficit from $q^a$ is at most
$kD_{a-1}$, where $k\ge1$ is the sum of the valuations of the chosen coefficients.
Multinomial coefficients can only increase the valuation. Hence such a term lies strictly
above the proposed face of slope $D_a$.

A monomial $\pi c_{a-1,j}X^j$ from the first summand has valuation $1+k$, where
$k=v(c_{a-1,j})$, and deficit
\[
 q^a-j=D_a+(q^{a-1}-j)\le D_a+kD_{a-1}.
\]
This is strictly smaller than $(1+k)D_a$ when $k>0$. When $k=0$ the monomial is precisely
$\pi X^{q^{a-1}}$. All other contributions at that exponent have valuation at least two, so
they cannot cancel its residue after division by $\pi$. Adding coefficients elsewhere can
only increase their valuations. The congruence of exponents follows by the same induction,
since a product of $q$ monomials with exponents congruent to $1$ has exponent congruent to
$q\equiv1$ modulo $q-1$.
\end{proof}

In particular
\begin{equation}\label{lub:eq:modsquare}
 \phi_a(X)\equiv X^{q^a}+\pi X^{q^{a-1}}\pmod{\pi^2}\qquad(a\ge1),
\end{equation}
which can also be seen directly: $\phi_{a-1}\equiv X^{q^{a-1}}\bmod\pi$, and replacing
$\phi_{a-1}$ by this monomial inside the $q$th power is valid modulo $\pi^2$ because
$\pi q\in\pi^2\OO$.

\begin{theorem}[Sharp pole bound and carry criterion; \stU]\label{lub:thm:pole}
Fix $b\in B$, $d\ge2$, and let $h_*=\max_{c\ne b}h_c$ be the deepest contact.
\begin{enumerate}[label=\textup{(\alph*)}]
\item If $h_*=0$, then $t_b(r)=\omega=d-1$ for every $r\ge1$.
\item If $h_*\ge1$, put
\[
 D=D_{h_*}=(q-1)q^{h_*-1},\qquad m=\nu_{h_*}=\#\set{c\ne b:\ h_c=h_*}.
\]
Then for every $r\ge1$
\begin{equation}\label{lub:eq:polebound}
 \omega+\max(0,r-m)\,D\ \le\ t_b(r)\ \le\ \omega+(r-1)D,
\end{equation}
and
\begin{equation}\label{lub:eq:carry}
 t_b(r)=\omega+(r-1)D\quad\Longleftrightarrow\quad
 p\nmid\binom{m+r-2}{r-1},
\end{equation}
equivalently, if and only if adding $m-1$ and $r-1$ in base $p$ produces no carry.
\end{enumerate}
In particular $t_b(1)=\omega$ in all cases.
\end{theorem}
\begin{proof}
By Lemma~\ref{lub:lem:normal} we compute with $P^{(b)}=\prod_{c\ne b}\phi_{h_c}$ in the
coordinate \eqref{lub:eq:coordinate}. If $h_*=0$, every factor is $X$ and
$P^{(b)}=X^{d-1}$, which proves (a).

Let $h_*\ge1$. For $a\ge1$ write $\phi_a=X^{q^a}(1+G_a)$, where $G_a$ is a polynomial in
$X^{-1}$ without constant term, all of whose coefficients are divisible by $\pi$; put
$G_0=0$. Modulo $\pi^r$,
\[
 \frac1{P^{(b)}}=X^{-\omega}\prod_{c\ne b}(1+G_{h_c})^{-1},\qquad
 (1+G_a)^{-1}=\sum_{k=0}^{r-1}(-G_a)^k .
\]
By Lemma~\ref{lub:lem:face}, every monomial $cX^{-\delta}$ of $G_a$ satisfies
$\delta\le D_a\,v(c)$. Since $D_a<D$ for $a<h_*$ and $v(c)\ge1$, all monomials of $G_a$
with $a<h_*$ satisfy the strict inequality $\delta<D\,v(c)$. For $a=h_*$ the same strict
inequality holds for every monomial of $G_{h_*}$ except one, $c_*X^{-D}$ with
$c_*\equiv\pi\pmod{\pi^2}$.

Expand the product. Each term is an integer multiple of a product of monomials of the
$G_{h_c}$, say $k$ copies of $c_*X^{-D}$ and further monomials $c_iX^{-\delta_i}$ of the
strict kind. The resulting monomial $cX^{-\delta}$ has $\delta\le D\,v(c)$, with strict
inequality as soon as one strict monomial occurs. A term that is nonzero modulo $\pi^r$ has
$v(c)\le r-1$, so the inverse has no nonzero exponent below $-\omega-(r-1)D$. This is the
upper bound in \eqref{lub:eq:polebound}. A term at exactly that exponent which survives
modulo $\pi^r$ contains no strict monomial and has $k=r-1$. The coefficient of
$X^{-\omega-(r-1)D}$ is therefore the coefficient supplied by
$X^{-\omega}(1+c_*X^{-D})^{-m}$, namely
\begin{equation}\label{lub:eq:topcoefficient}
 (-c_*)^{r-1}\binom{m+r-2}{r-1}\equiv(-\pi)^{r-1}\binom{m+r-2}{r-1}\pmod{\pi^r}.
\end{equation}
It is nonzero in $A_r$ exactly when the integer binomial coefficient is prime to $p$. This
proves \eqref{lub:eq:carry}. The carry formulation is Kummer's criterion: by
Legendre's formula $v_p(n!)=(n-s_p(n))/(p-1)$, with $s_p$ the sum of base-$p$ digits,
\[
 v_p\binom{x+y}{x}=\frac{s_p(x)+s_p(y)-s_p(x+y)}{p-1},
\]
and this vanishes exactly when the addition of $x$ and $y$ has no carry.

For the lower bound put $a_k=\binom{m+k-1}{k}\bmod p$. The generating series of the $a_k$
is $(1-z)^{-m}\in\F_p[[z]]$. The recurrence obtained from multiplication by $(1-z)^m$,
whose constant and leading coefficients are nonzero, forbids $m$ consecutive zeros among
the $a_k$: such a block would force all later coefficients to vanish, while $(1-z)^{-m}$ is
not a polynomial. Since $a_0=1$, some $j$ with $\max(0,r-m)\le j\le r-1$ has $a_j\ne0$.
The pole order is weakly increasing in the precision, because integrality modulo $\pi^r$
implies integrality modulo $\pi^{j+1}$. With \eqref{lub:eq:carry} at precision $j+1$ this
gives $t_b(r)\ge t_b(j+1)=\omega+jD$.
\end{proof}

\begin{corollary}[Residue field; unique deepest contact; \stU]\label{lub:cor:unique}
\begin{enumerate}[label=\textup{(\alph*)}]
\item Let $r=1$, so that $A_1=k$. The observations determine the jet of every
$H\in k[[X,Y]]$ if and only if $n_b\ge1+\omega_b$ for every $b\in B$. Here
$\omega_b=\sum_{c\ne b}q^{v(b-c)}$ is the sum of the contact orders of the branch $b$
with the other branches. The least common precision is $1+\max_b\omega_b$.
\item If the deepest contact of $b$ is unique \textup{(}$m=1$\textup{)}, then
$t_b(r)=\omega+(r-1)D$ for every $r$. In particular, for exactly two branches with
$h=v(b-c)\ge1$,
\[
 t_b(r)=q^h+(r-1)(q-1)q^{h-1},
\]
and both branches have the same least precision $1+t_b(r)$.
\item If $h_*\ge1$, then $t_b(r)-Dr$ is bounded and $\lim_{r\to\infty}t_b(r)/r=D$. Thus the
necessary branch depth grows with $r$ as soon as two directions are congruent modulo $\pi$.
\end{enumerate}
\end{corollary}
\begin{proof}
Combine Proposition~\ref{lub:prop:threshold} with Theorem~\ref{lub:thm:pole}; in (b) the binomial coefficient
in \eqref{lub:eq:carry} equals $1$.
\end{proof}

\begin{remark}[The weaker slope $q^{h_*}-1$]\label{lub:rem:older}
Allowing each power of $\pi$ in the geometric expansion of Lemma~\ref{lub:lem:contact} to add at
most $q^{h_*}-1$ to the pole order gives the cruder estimate
$t_b(r)\le\omega+(r-1)(q^{h_*}-1)$. Since $D\le q^{h_*}-1$ with equality exactly for
$h_*=1$, Theorem~\ref{lub:thm:pole} shows that this estimate is attained in precisely the
following cases: always for $r=1$; always for $h_*=0$; for $r>1$ and $h_*=1$ exactly when
$(m-1)+(r-1)$ has no carry in base $p$; and never for $r>1$ and $h_*\ge2$.
\end{remark}

\begin{remark}[Frequency of equality]\label{lub:rem:frequency}
If the base-$p$ digits of $m-1$ are $a_0,\dots,a_{\ell-1}$, the upper endpoint in
\eqref{lub:eq:polebound} is attained for a proportion $\prod_{i=0}^{\ell-1}(p-a_i)/p^{\ell}$
of the precisions $r$ in every block of $p^{\ell}$ consecutive values, by independent digit
choices in $r-1$.
\end{remark}

\begin{remark}[A classical special case]\label{lub:rem:classical}
For $K=\Qp$, $\pi=p$ and two directions with contact depth $h\ge1$, the multiplicative
coordinate makes the denominator associated to $(1+X)^{p^h}-1$. By
Lemma~\ref{lub:lem:normal}(iii) the condition $X^n/\Delta_b\in A_r[[X]]$ says that $X^n=0$ in
the quotient by this polynomial, which is $(\Z/p^r)[T]/(T^{p^h}-1)$ via $X=T-1$. The value
$p^h+(r-1)(p-1)p^{h-1}$ of Corollary~\ref{lub:cor:unique}(b) agrees with the classical nilpotency
index of the augmentation ideal of this cyclic group ring, which follows from results of
Weisman and Wilson as presented in \cite[Theorem~4.7]{ClarkSchauz2022Functional}. More than
two branches give products of division polynomials in the same variable, with
multiplicities, and allow the cancellation in \eqref{lub:eq:carry}.
\end{remark}

\subsection{Exact pole orders below saturation}\label{lub:ssec:exact}
\begingroup\raggedright\sources{grh_prime_patterns_and_branch_precision; grh_transfers_lubin_tate_and_graph_towers; grh_method_transfers; grh_method_transfers_orthogonal_jets_2026-11-26}\endgroup

When the carry criterion \eqref{lub:eq:carry} fails, the shallower contacts can affect the
exact pole order. The next theorem computes it at every precision from one finite ring
that does not depend on $r$.

Keep the notation of Theorem~\ref{lub:thm:pole} with $h_*\ge1$ and put $\kappa=q^{h_*-1}$, so that
$D=(q-1)\kappa$. By the last assertion of Lemma~\ref{lub:lem:face},
\begin{equation}\label{lub:eq:F}
 X^{-\omega}\prod_{c\ne b}\phi_{h_c}(X)=F\bigl(X^{-(q-1)}\bigr),\qquad F\in\OO[z],\quad F(0)=1 .
\end{equation}
In particular the defect $\omega+(r-1)D-t_b(r)$ is always divisible by $q-1$.

\begin{theorem}[The exact defect at every precision; \stU]\label{lub:thm:defect}
Let $h_*\ge1$. There is a bounded, eventually periodic sequence of integers $\eta_r$ with
\[
 t_b(r)=\omega+(q-1)\bigl(\kappa(r-1)-\eta_r\bigr),\qquad 0\le\eta_r\le(m-1)\kappa .
\]
Its exact values, including the initial ones, are computed as follows. Adjoin $\rho$ with
$\rho^{\kappa}=\pi$, put $N=m\kappa$ and $S=\OO[\rho]/(\rho^N)$, with $v_\rho(\rho)=1$ and
the convention that the zero class of $S$ has valuation $+\infty$. Set
\[
 G(z)=F(z/\rho),\qquad \frac1{G(z)}=\sum_{j\ge0}\gamma_jz^j\ \text{ in }S[[z]],\qquad
 \gamma_j=0\ (j<0).
\]
Then $G$ has coefficients in $\OO[\rho]$, and $\eta_r$ is the least integer
$i\in[0,(m-1)\kappa]$ satisfying
\begin{equation}\label{lub:eq:survival}
 v_\rho\bigl(\gamma_{\kappa(r-1)-i}\bigr)<\kappa+i .
\end{equation}
Moreover, with
\[
 N_1=mN,\qquad b_0=\ceil{m/e},\qquad
 \Theta=2p^{\,b_0+\floor{\log_p\max\set{1,N_1-1}}},
\]
one has $\eta_{r+\Theta}=\eta_r$ for all $r$ with $\kappa(r-m)\ge\deg C$, where
$C\in S[z]$ is the polynomial \eqref{lub:eq:C} below. The period $\Theta$ is explicit and
possibly not minimal.
\end{theorem}
\begin{proof}
The ring $\OO[\rho]$ is a discrete valuation ring with uniformizer $\rho$ and
$v_\rho(\pi)=\kappa$. By the degree allowance established in the proof of
Theorem~\ref{lub:thm:pole}, a monomial $cz^j$ of $F$ corresponds to $cX^{-(q-1)j}$ with
$(q-1)j\le D\,v(c)$, that is $\kappa\,v(c)\ge j$. So $G$ is integral. On reducing $G$
modulo $\rho$, only the terms with $\kappa\,v(c)=j$ survive; these are the constant term
and the products of the monomials $c_*z^{\kappa}$ from the $m$ deepest factors, and
$c_*\rho^{-\kappa}=c_*/\pi\equiv1$. Thus
\begin{equation}\label{lub:eq:Gmod}
 G(z)\equiv(1+z^{\kappa})^m\pmod\rho .
\end{equation}
Write $F^{-1}=\sum_j\theta_jz^j$ with $\theta_j\in\OO$. Then $\gamma_j$ is the class of
$\theta_j\rho^{-j}$, and by \eqref{lub:eq:F} and Lemma~\ref{lub:lem:normal}
\[
 t_b(r)=\omega+(q-1)\max\set{j:\ v(\theta_j)<r}.
\]
At $j=\kappa(r-1)-i\ge0$ the condition $v(\theta_j)<r$ is equivalent to
$v_\rho(\theta_j\rho^{-j})<\kappa+i$, since $v(\theta_j)$ is an integer. All these cutoffs
are at most $N$, so the test can be made in $S$. By \eqref{lub:eq:polebound} the largest
surviving index $j$ lies in the range $\kappa(r-1)-(m-1)\kappa\le j\le\kappa(r-1)$. This
proves the formula for $t_b(r)$, the bounds on $\eta_r$ and the test \eqref{lub:eq:survival}.

For the period put $E=G-(1+z^{\kappa})^m\in S[z]$ and
\begin{equation}\label{lub:eq:C}
 C(z)=\sum_{k=0}^{N-1}(-E)^k(1+z^{\kappa})^{m(N-1-k)} .
\end{equation}
By \eqref{lub:eq:Gmod} all coefficients of $E$ lie in $\rho S$, so $E^N=0$ and
$G^{-1}=C(z)/(1+z^{\kappa})^{N_1}$ in $S[[z]]$. For $j\ge\deg C$, the coefficient
$\gamma_j$ is a finite $S$-linear combination of the shifted sequences
$(-1)^u\binom{u+N_1-1}{N_1-1}$, indexed by $u=(j-k)/\kappa$ for the monomials
$z^k$ of $C$ with $k\equiv j\bmod\kappa$. Now $p^{b_0}=0$ in $S$, because
$v_\rho(p)=\kappa e$. For $1\le l\le N_1-1$ the identity
$l\binom{\Theta}{l}=\Theta\binom{\Theta-1}{l-1}$ gives $p^{b_0}\mid\binom{\Theta}{l}$.
Vandermonde's identity then shows that the displayed signed binomial sequence has period
dividing $\Theta$ in $S$; this includes $N_1=1$, where the evenness of $\Theta$
suffices. Hence $\gamma_{j+\kappa\Theta}=\gamma_j$ for $j\ge\deg C$, and
$\eta_{r+\Theta}=\eta_r$ whenever $\kappa(r-1)-(m-1)\kappa\ge\deg C$.
\end{proof}

\begin{remark}[Computation over the original local ring]\label{lub:rem:original}
The auxiliary root $\rho$ can be removed. Define $k_j=\ceil{j/\kappa}$,
$c_i=[z^i]F/\pi^{k_i}$ and $\beta_j=\theta_j/\pi^{k_j}$. Dividing the recurrence
$[z^j](FF^{-1})=0$ by $\pi^{k_j}$ gives $\beta_0=1$ and, for $j>0$,
\[
 \beta_j=-\sum_{i=1}^{\min(j,\deg F)}c_i\,\beta_{j-i}\,\pi^{k_i+k_{j-i}-k_j}\pmod{\pi^m}.
\]
The exponents of $\pi$ in this expression are zero or one. For $j=\kappa(r-1)-i\ge0$ the
survival test \eqref{lub:eq:survival} becomes
$\beta_j\not\equiv0\bmod\pi^{\floor{i/\kappa}+1}$. The recurrence computes every precision
in the fixed ring $\OO/\pi^m$.
\end{remark}

In two situations the defect has a closed form.

\begin{proposition}[All precisions when every contact has depth at most one; \stU]\label{lub:prop:depthone}
Suppose $h_*=1$, so that every $h_c\in\set{0,1}$ and $m\ge1$ counts the contacts of depth
one. Set
\[
 K_m(r)=\max\Bigl\{0\le k<r:\ k+e\,v_p\binom{m+k-1}{k}<r\Bigr\}.
\]
Then
\begin{equation}\label{lub:eq:depthone}
 t_b(r)=\omega+(q-1)K_m(r)\qquad(r\ge1).
\end{equation}
\end{proposition}
\begin{proof}
In the coordinate \eqref{lub:eq:coordinate} the unit-free denominator is exactly
$X^{\omega}(1+\pi X^{-(q-1)})^m$. Its reciprocal modulo $\pi^r$ is
\[
 X^{-\omega}\sum_{k=0}^{r-1}(-1)^k\binom{m+k-1}{k}\pi^kX^{-k(q-1)} .
\]
Distinct $k$ give distinct exponents. The coefficient at index $k$ has valuation
$k+e\,v_p\binom{m+k-1}{k}$, so it survives exactly under the condition in the maximum; the
candidate set contains $k=0$. No cancellation between different $k$ is possible.
\end{proof}

\begin{corollary}[An explicit period at depth one; \stU]\label{lub:cor:depthoneperiod}
In Proposition~\ref{lub:prop:depthone} let $m\ge2$, let $p^{\sigma}$ be the least power of $p$ with
$p^{\sigma}\ge m$, and put
\[
 \ell=\floor{(p^{\sigma}-1)/e},\qquad \vartheta=\ell+v_p\bigl((m-1)!\bigr)+1 .
\]
Then $K_m(r)\ge\max(0,r-p^{\sigma})$ for every $r\ge1$, and for $r\ge p^{\sigma}$
\[
 K_m(r+p^{\vartheta})=K_m(r)+p^{\vartheta},\qquad\text{so that}\qquad
 t_b(r+p^{\vartheta})=t_b(r)+(q-1)p^{\vartheta}.
\]
For $m=1$ one has $K_1(r)=r-1$ from the start.
\end{corollary}
\begin{proof}
Every $k$ divisible by $p^{\sigma}$ has no carry when added to $m-1<p^{\sigma}$, so its
binomial coefficient is a $p$-adic unit. The largest such $k\le r-1$ proves the first
assertion. For $r\ge p^{\sigma}$ write a potential maximum as $k=r-1-\delta$ with
$0\le\delta<p^{\sigma}$. With $n=m-1$ it is
admissible exactly when
\[
 e\,v_p\binom{k+n}{n}\le\delta .
\]
All thresholds for the binomial valuation are at most $\ell$. Since
$\binom{k+n}{n}=\prod_{i=1}^{n}(k+i)/n!$, these comparisons are determined by
$k\bmod p^{\vartheta}$: if the valuation of a numerator factor reaches $\vartheta$, the
binomial valuation already exceeds every relevant threshold; below that, the valuation of
each factor is determined by its residue modulo $p^{\vartheta}$. Replacing $r$ by
$r+p^{\vartheta}$ leaves the set of admissible $\delta$ unchanged, hence also its smallest
element, which gives the largest $k$.
\end{proof}

The period $p^{\vartheta}$ is a safe explicit period and is not claimed to be minimal;
the finite maximum in \eqref{lub:eq:depthone} is usually the simpler algorithm.

\begin{proposition}[Second precision at arbitrary contact depths; \stU]\label{lub:prop:second}
For every contact multiset,
\[
 t_b(2)=\omega+\max\bigl(\set{0}\cup\set{D_j:\ j\ge1,\ p\nmid\nu_j}\bigr).
\]
\end{proposition}
\begin{proof}
By \eqref{lub:eq:modsquare}, the reciprocal of $P^{(b)}$ modulo $\pi^2$ is
\[
 X^{-\omega}\Bigl(1-\pi\sum_{j\ge1}\nu_jX^{-D_j}\Bigr).
\]
The exponents $D_j$ are distinct, and the term indexed by $j$ survives modulo $\pi^2$
exactly when $\nu_j$ is prime to $p$. The largest surviving level determines the pole.
\end{proof}

\begin{example}\label{lub:ex:poles}
\begin{enumerate}[label=\textup{(\arabic*)}]
\item Take $K=\Q_3$, $\pi=3$ and $B=\set{0,3,6}$. Every pair of directions has contact
depth one, so at each branch $\omega=6$ and $m=2$, and \eqref{lub:eq:depthone} gives the
following table. The estimate of Remark~\ref{lub:rem:older} is listed for comparison.
\begin{center}\small
\begin{tabular}{l|rrrrrrrrr}
\toprule
$r$ &1&2&3&4&5&6&7&8&9\\
\midrule
$K_2(r)$ &0&1&1&3&4&4&6&7&7\\
exact precision $1+t_b(r)$ &7&9&9&13&15&15&19&21&21\\
$1+\omega+(r-1)(q^{h_*}-1)$ &7&9&11&13&15&17&19&21&23\\
\bottomrule
\end{tabular}
\end{center}
At $r=3$ the coefficient at $k=2$ is $3^2\binom32=27\equiv0$, so observations through
degree eight suffice and through degree seven they do not. The plateaus are arithmetic
cancellations, not rounding effects.
\item Take $K=\Q_2$, $\pi=2$, $B=\set{0,2,6}$ and $b=0$. The contact multiset is $(1,1)$,
the denominator is associated to $X^2(2+X)^2$, and the reciprocal terms have coefficients
$(-2)^k(k+1)$. Hence $t_0(r)=4+(r-1)$ for odd $r$ and $t_0(r)=4+(r-2)$ for even $r$; at
precisions one and two the same five coefficients suffice at this branch.
\item Take $K=\Q_2$, $B=\set{0,4,12}$ and $b=0$. Here $\omega=8$, $D=2$ and $m=2$. The pole
order is $8$ at $r=2$, where $\phi_2^2\equiv X^8\bmod4$, and $12$ at $r=3$, where the top
coefficient survives; the observation precisions are $9$ and $13$. For $B=\set{0,4}$ and
$r=2$ one finds $t=6$: in the multiplicative coordinate $[4](X)=(1+X)^4-1$ reduces to
$X^4+2X^2$ modulo $4$, with reciprocal $X^{-4}-2X^{-6}$.
\item For $K=\Q_3$, two directions that differ by $9$ have $\omega=9$ and $D=6$, and their
common cutoff is $10+6(r-1)$ at every $r$.
\item Let $p=2$ and $q=4$, and compare the contact multisets $(0,2,2)$ and $(1,2,2)$. Both
have $D=12$ and $m=2$, and neither reaches $\omega+D$ at $r=2$. By
Proposition~\ref{lub:prop:second} the defect is $12$ in the first case. In the second the depth-one
factor leaves a surviving pole increase of $3$, and the defect is $9$. Thus shallower
contacts can determine the exact answer already at the first cancellation.
\end{enumerate}
\end{example}

\begin{remark}[Finite checks]\label{lub:rem:checks}
The formulas of this subsection, of the preceding one, and of
Theorem~\ref{lub:thm:directions} below were compared with exact computations in the
delivered files:
\begin{itemize}[nosep]
\raggedright
\item $210$ inverse products for $(p,q)=(2,2),(2,4),(3,3),(3,9)$ and $r\le6$:
\artifact{grh_method_transfers_orthogonal_jets_2026-11-26/checks/jet_checks.json};
\item $1400$ monic-remainder computations and $640$ comparisons with the literal branches
$(1+X)^b-1$:
\artifact{grh_transfers_lubin_tate_and_graph_towers/checks/branch_check_output.txt};
\item observation matrices over $\Z/p^r$ for small sets such as $\set{0,2,6}$ over $\Z_2$
and $\set{0,3,6}$ over $\Z_3$:
\artifact{grh_transfers_lubin_tate_and_graph_towers/checks/jet-matrix-output.txt};
\item $1088$ cases of the bounds \eqref{lub:eq:polebound} and of saturation:
\artifact{grh_prime_patterns_and_branch_precision/results/branch_checks.txt};
\item $6128$ finite residue-ball designs for the minimax formula \eqref{lub:eq:minimax}:
\artifact{grh_transfers_lubin_tate_and_graph_towers/checks/design_checks.txt}.
\end{itemize}
These computations use unramified coefficient rings; the proofs, including the factor $e$,
cover ramified fields.
\end{remark}

\subsection{Optimal directions and a multiplication rule}\label{lub:ssec:directions}
\begingroup\raggedright\sources{grh_transfers_cycle_signings_and_lubin_tate; grh_transfers_lubin_tate_and_graph_towers; grh_prime_patterns_and_branch_precision}\endgroup

In \cite[Theorem~4.6]{BergerSprang2025Integer} the directions must lie in $\pi\OO$ in order
to preserve the kernel of the trace operator $\psi$. This motivates optimizing the common
precision over all sets of $d$ directions inside a prescribed residue ball.

\begin{theorem}[Optimal directions at all precisions in every residue ball; \stU]\label{lub:thm:directions}
Suppose $d\ge2$ directions may be chosen in a prescribed coset of $\pi^a\OO$, $a\ge0$. The
smallest possible common precision is
\begin{equation}\label{lub:eq:minimax}
 1+\tau_r\Bigl(X^{d-1}\prod_{j\ge1}\Phi_j^{\,u_j}\Bigr),\qquad
 u_j=\begin{cases} d-1,& j\le a,\\[2pt] \floor{(d-1)/q^{\,j-a}},& j>a.\end{cases}
\end{equation}
Equivalently, it is one plus the pole order of the reciprocal of
$\prod_{i=1}^{d-1}\phi_{H_i}$, where $H_i=a+\max\set{h:\ q^h\mid i}$. One set of directions
is optimal simultaneously at every precision $r$: choose representatives
$x_0,\dots,x_{q-1}$ of the residue field, map the base-$q$ digits of $0,\dots,d-1$ to these
representatives in a $\pi$-adic expansion, scale by $\pi^a$, and translate into the
required coset.
\end{theorem}
\begin{proof}
Translating and scaling, the directions are $\xi+\pi^ay$ with a fixed $\xi\in\OO$ and $y$
running over a set $B_0\subset\OO$ of $d$ elements, and contact depths in $B$ are those in
$B_0$ increased by $a$. Hence $m_j(b)=d-1$ for
$j\le a$ and every $b$. Starting with $B_0$, successively follow a most populous child in
the tree of residue balls. By induction its occupancy after $h$ splits is at least
$\ceil{d/q^h}$. The nested path selects some $b$, since the finite set is nonempty at every
depth and eventually the ball contains a single direction. For this $b$ the number of
other directions with $v(b-c)\ge a+h$ is at least
$\ceil{d/q^h}-1=\floor{(d-1)/q^h}$, that is $m_j(b)\ge u_j$ for all $j$. By
\eqref{lub:eq:denominator} the polynomial in \eqref{lub:eq:minimax} divides $P^{(b)}$, and
Lemma~\ref{lub:lem:normal}(iii) gives the lower bound.

In the digit construction, two directions are congruent modulo $\pi^{a+h}$ exactly when
their indices are congruent modulo $q^h$. Every residue class at depth $h$ therefore
contains at most $\ceil{d/q^h}$ directions, so $m_j(b)\le u_j$ for every $b$ and $j$. Now
$P^{(b)}$ divides the polynomial in \eqref{lub:eq:minimax}, which gives the matching upper
bound at every $r$. The second description holds because the number of
$i\in\set{1,\dots,d-1}$ with $H_i\ge j$ is $u_j$.
\end{proof}

\begin{corollary}[Closed forms for the optimum; \stU]\label{lub:cor:optimum}
\begin{enumerate}[label=\textup{(\alph*)}]
\item Over the residue field, that is for $r=1$, the least common precision for $d$
directions in a coset of $\pi^a\OO$ equals $1+q^aW_q(d)$, where
\[
 W_q(d)=d-1+(q-1)\sum_{h\ge1}q^{h-1}\floor*{\frac{d-1}{q^h}} .
\]
For $d=q^t$, $t\ge1$, this is $1+q^a\,t(q-1)q^{t-1}$.
\item For $d=q^k+1$ with $k\ge1$ the extremal contact list has a unique deepest contact,
and the optimal common precision is, for every $r\ge1$,
\[
 1+q^a\bigl(q^k+k(q-1)q^{k-1}\bigr)+(r-1)(q-1)q^{a+k-1}.
\]
\end{enumerate}
\end{corollary}
\begin{proof}
(a) Modulo $\pi$ the polynomial in \eqref{lub:eq:minimax} is a power of $X$, so its
$\tau_1$ is its degree $d-1+\sum_jD_ju_j=q^aW_q(d)$. When $d=q^t$, precisely $(q-1)q^{t-h-1}$ other
directions have contact depth $a+h$ with a given one, for $0\le h<t$.
(b) The polynomial in \eqref{lub:eq:minimax} is $P^{(b)}$ for the direction of index $0$
in the digit construction. Counting the multiples of $q^h$ gives
$\omega=q^a(q^k+k(q-1)q^{k-1})$; the deepest level $a+k$ has multiplicity $u_{a+k}=1$, and
Corollary~\ref{lub:cor:unique}(b) applies.
\end{proof}

\begin{example}\label{lub:ex:eight}
For eight directions over $\Q_2$ the profiles $(u_j)$ in \eqref{lub:eq:minimax} are $(3,1)$
without restriction and $(7,3,1)$ when all eight directions must be even; in both the
deepest level has multiplicity one. The optimal common precision is $2r+11$ in the first
case and $4r+21$ in the second. At $r=1$ these are $13$ and $25$: inside $\pi\OO$, all
bivariate coefficients of total degree at most seven are recovered from branch
coefficients through degree $24$. For $q=2$, five directions constrained to one coset of
$\pi\OO$ need $17+4(r-1)$ coefficients per branch. These residue-class optimizations are
related to Bhargava's $P$-orderings \cite{Bhargava1997Porderings}.
\end{example}

Berger and Sprang use the case $r=1$ to multiply coefficient polynomials
\cite[Theorem~4.6]{BergerSprang2025Integer}. For $f\in\OO[[X]]$ and $x\in\OO$ write
$c_{f,j}(x)=[X^j]f([x](X))$. The interpolation formula \eqref{lub:eq:lagrange} gives an
explicit row cutoff for this procedure at any residue precision.

\begin{proposition}[A finite-precision multiplication rule; \stU]\label{lub:prop:products}
Let $\psi$ be the Lubin--Tate trace operator of \cite{BergerSprang2025Integer}, let
$f\in\ker\psi$ and $g\in\OO[[X]]$, and let $i,\ell\ge0$. Set $d=i+\ell+1$, choose
$B\subset\pi\OO$ of size $d$ by Theorem~\ref{lub:thm:directions}, and define
$f_b(X)=f(X)\,g([b](X))$ for $b\in B$. Then $f_b\in\ker\psi$, and with the constants
\[
 \mu_{bj}=[X^{i-j}Y^{\ell}]\,\frac{P_b(X,Y)}{\Delta_b(X)}\in A_r
\]
the following congruence of integer-valued functions holds for every $x\in\OO$:
\begin{equation}\label{lub:eq:products}
 c_{f,i}(x)\,c_{g,\ell}(x)\equiv\sum_{b\in B}\sum_{j=0}^{t_b(r)}\mu_{bj}\,c_{f_b,j}(x)\pmod{\pi^r}.
\end{equation}
\end{proposition}
\begin{proof}
The membership $f_b\in\ker\psi$ is \cite[Lemma~4.5]{BergerSprang2025Integer}. Apply
\eqref{lub:eq:lagrange} to
\[
 H(X,Y)=f([x](X))\,g([x](Y)),
\]
whose coefficient at $X^iY^{\ell}$, of total degree $<d$, is $c_{f,i}(x)c_{g,\ell}(x)$.
The branch values are $H(X,u_b)=f_b([x](X))$, since the endomorphisms commute. Moreover
$[Y^{\ell}]P_b$ lies in $X^{d-1-\ell}A_r[[X]]$ and $1/\Delta_b$ lies in
$X^{-t_b(r)}A_r[[X]]$, so $\mu_{bj}=0$ for $j>t_b(r)$. This gives the displayed cutoff.
\end{proof}

Equation \eqref{lub:eq:products} supplies explicit weights and an optimal common precision
for the full jet used by this procedure. A single prescribed product can require fewer
coefficients.

\subsection{A simultaneous congruence for weighted Fleck sums}\label{lub:ssec:fleck}
\begingroup\raggedright\sources{grh_prime_patterns_and_branch_precision}\endgroup

The multiple-branch calculation has an arithmetic interpretation involving only ordinary
integers. Fix a prime $p$, a modulus $M=p^h$ with $h\ge1$, and a positive integer $m$. For
$0\le u<M$ and $0\le j<m$ define
\[
 S_{j,u}(n)=\sum_{\substack{0\le k\le n\\ k\equiv u\bmod M}}
 (-1)^{n-k}\binom nk\binom{(k-u)/M}{j}.
\]
The lower argument $(k-u)/M$ is a nonnegative integer at every summand. Sums of this type
are, up to the displayed sign, the cyclotomic $\psi$-coefficients before normalization, for
which Sun and Wan prove Lucas-type congruences \cite{SunWan2005Lucas}. Their integrality
statement gives
$v_p(S_{j,u}(n))\ge\floor{(n-M/p-jM)/(M-M/p)}$; at $j=m-1$ it supplies the sufficient
common threshold $U_r$ below. The following result determines the least simultaneous
threshold and identifies the obstruction one step earlier.

\begin{theorem}[Simultaneous congruence for weighted Fleck sums; \stU]\label{lub:thm:fleck}
For $r\ge1$ let $T_r$ be the inverse pole order \eqref{lub:eq:pole} for $K=\Qp$, $\pi=p$
and the contact multiset consisting of $m$ copies of $h$; thus $\omega=mM$, $D=M-M/p$, and
with $U_r=mM+(r-1)D$
\[
 U_r-(m-1)D\le T_r\le U_r,\qquad
 T_r=U_r\iff\text{$r-1$ and $m-1$ add without a base-$p$ carry.}
\]
Then:
\begin{enumerate}[label=\textup{(\alph*)}]
\item $p^r\mid S_{j,u}(n)$ for every $0\le j<m$ and $0\le u<M$ if and only if $n\ge T_r$.
\item For every $r$, the values $S_{m-1,u}(T_r-1)$, $0\le u<M$, are all congruent modulo
$p^r$ to one nonzero class $c$, and $S_{j,u}(T_r-1)\equiv0\pmod{p^r}$ for $j<m-1$.
\item When $T_r=U_r$, the common class is
\[
 c\equiv(-p)^{r-1}\binom{m+r-2}{r-1}\pmod{p^r}.
\]
\end{enumerate}
\end{theorem}
\begin{proof}
The displayed bounds are Theorem~\ref{lub:thm:pole}. Put $\Xi=Y^M-1$ and $\Pi=\Xi^m$ in
$(\Z/p^r)[Y]$. In the multiplicative Lubin--Tate coordinate $X=Y-1$ one has
$[p^h](X)=(1+X)^M-1$, so by Lemma~\ref{lub:lem:normal} the number $T_r$ is the nilpotency index
of $Y-1$ in $(\Z/p^r)[Y]/(\Pi)$. The polynomials $Y^u\Xi^j$ with $0\le u<M$ and $0\le j<m$
form a basis of this quotient over $\Z/p^r$: their leading monomials run through the
degrees $0$ to $mM-1$ with coefficient one. Writing $Y^{u+Ml}=Y^u(1+\Xi)^l$ in the binomial
expansion of $(Y-1)^n$ identifies the coordinates of $(Y-1)^n$ in this basis with the
$S_{j,u}(n)$. This proves (a).

Let $Q(Y)$ be the polynomial remainder of $(Y-1)^{T_r-1}$ on division by $\Pi$. Then
$Q\ne0$ and $(Y-1)Q=c\,\Pi$ for some scalar $c\in\Z/p^r$, by the degree bound and the
vanishing at $T_r$. Multiplication by the monic polynomial $Y-1$ is injective in
$(\Z/p^r)[Y]$, whence
\[
 Q=c\,\Xi^{m-1}(1+Y+\cdots+Y^{M-1}),\qquad c\ne0 .
\]
This is the full vector of residues claimed in (b).

For (c), put $\Pi_X=((1+X)^M-1)^m$ and $F_*(z)=z^{mM}\Pi_X(1/z)$. In the variable $X$ the
coefficient in degree $mM-1$ of the remainder of $X^{T_r-1}$ is still $c$. Monic division
at infinity identifies it with $[z^{T_r-mM}]F_*^{-1}$, by comparing the coefficients of
$X^{-1}$ in $X^{T_r-1}/\Pi_X$. In the reverse polynomial of a single factor the nonconstant
coefficients are $\binom Mi$ in deficit $M-i$, $1\le i<M$. They have valuations
$h-v_p(i)$, since $\binom Mi=(M/i)\binom{M-1}{i-1}$ and $\binom{M-1}{i-1}$ is prime to $p$.
The unique coefficient whose deficit equals $D$ times its valuation is at $i=M/p$; for all
others the deficit is strictly smaller than $D$ times the valuation. The coefficient at
$i=M/p$, divided by $p$, is
$\binom{M-1}{M/p-1}\equiv1\pmod p$. So $F_*$ has the same leading face as $(1+pz^D)^m$, and
the top-degree computation in the proof of Theorem~\ref{lub:thm:pole}, applied to
$(1+pz^D)^{-m}$, gives the residue \eqref{lub:eq:topcoefficient} with $\pi=p$.
\end{proof}

\begin{example}\label{lub:ex:fleck}
\sloppy
Take $p=2$, $h=2$, $m=2$ and $r=2$. The upper endpoint is $U_2=10$ and the exact threshold
is $T_2=8$: indeed $(Y^4-1)^2\equiv(Y-1)^8\bmod4$ by the binomial theorem, and at $n=7$
direct expansion gives $S_{1,u}(7)\equiv1\bmod4$ for every $u$. These residues, and both
sides of the threshold for further parameters ($3153$ vectors of exact integer sums), are
recorded in \artifact{grh_prime_patterns_and_branch_precision/results/fleck_checks.txt}.
\end{example}

\subsection{The logarithmic coefficient window and the first coefficient depth}\label{lub:ssec:depth}
\begingroup\raggedright\sources{grh_transfer_advances_2026-11-26; grh_transfers_cycle_signings_and_lubin_tate}\endgroup

Let $p$ be a prime and $q$ a power of $p$, and let $L(X)\in\Qp[[X]]$ be the normalized
Lubin--Tate logarithm with
\[
 L(X)=X+O(X^2),\qquad L(pX+X^q)=pL(X).
\]
Here $v_p$ is normalized by $v_p(p)=1$, and $v_p(0)=+\infty$.

\begin{theorem}[Exact logarithmic coefficient window; \stU]\label{lub:thm:window}
For every positive integer $n$ put $N=v_p(n)$. For every $1\le j\le p^N$,
\begin{equation}\label{lub:eq:window}
 v_p\Bigl([X^{n+j(q-1)}]\,L(X)^n\Bigr)=N-v_p(j)-j .
\end{equation}
The coefficients of $L(X)^n$ outside the degrees congruent to $n$ modulo $q-1$ vanish.
Consequently the first nonintegral coefficient of $L(X)^n$ after the leading one occurs
exactly in degree
\[
 n+(q-1)J_p(N),\qquad J_p(N)=\min\set{j\ge1:\ j+v_p(j)>N}.
\]
\end{theorem}
\begin{proof}
Choose $\eta$ with $\eta^{q-1}=p$ and write
\[
 \Lambda(X)=\eta^{-1}L(\eta X)=X\sum_{j\ge0}\lambda_jX^{j(q-1)} .
\]
The functional equation becomes $\Lambda(pX+pX^q)=p\Lambda(X)$. Coefficient comparison
gives $\lambda_0=1$ and, for $j\ge1$,
\[
 \bigl(1-p^{j(q-1)}\bigr)\lambda_j=\sum_{i=0}^{j-1}\lambda_i\,p^{i(q-1)}\binom{1+i(q-1)}{j-i}.
\]
Every factor on the left is a unit. Induction proves $\lambda_j\in\Zp$,
\[
 \lambda_1=(1-p^{q-1})^{-1},\qquad \lambda_j\in p^{q-1}\Zp\quad(j\ge2).
\]
The support assertion follows either from the same recurrence or from the $\mu_{q-1}$
symmetry of the functional equation.

Let $C(Y)=\sum_j\lambda_jY^j$ and $C_0(Y)=C(Y)-1$. There is $\beta(Y)\in Y\Zp[[Y]]$ such
that
\[
 C_0(Y)=\lambda_1Y\bigl(1+p^{q-1}\beta(Y)\bigr).
\]
For every integer $i\ge1$, all nonconstant coefficients of $(1+p^{q-1}\beta)^i$ are
divisible by $p^{q-1+v_p(i)}$. Indeed $v_p\binom ik\ge v_p(i)-v_p(k)$ and
$(q-1)k-v_p(k)\ge q-1$ for $k\ge1$. It follows that for $j>i$,
\[
 v_p\bigl([Y^j]C_0(Y)^i\bigr)\ge q-1+v_p(i).
\]
Expand $C^n=\sum_i\binom niC_0^i$. For $1\le j\le p^N\le n$, the term $i=j$ contributes
$\binom nj\lambda_1^j$ to the coefficient of $Y^j$. Its valuation is $N-v_p(j)$: use
$\binom nj=(n/j)\binom{n-1}{j-1}$ and note that the last binomial coefficient is a $p$-adic
unit, since the lowest $N$ base-$p$ digits of $n-1$ are all $p-1$. Every term with $i<j$
has valuation at least $N+q-1$, by the preceding inequality and
$v_p\binom ni\ge N-v_p(i)$. Thus no cancellation is possible and
\[
 v_p\bigl([Y^j]C^n\bigr)=N-v_p(j).
\]
Since $[X^{n+j(q-1)}]\Lambda^n=p^j[X^{n+j(q-1)}]L^n$, equation \eqref{lub:eq:window}
follows. Finally $J_p(N)\le N+1\le p^N$, also when $N=0$, so the proven range includes the
first failure of integrality.
\end{proof}

\begin{remark}\label{lub:rem:Jp}
\sloppy
The case $p=q=2$ gives $L(X)=\log(1+X)$, and \eqref{lub:eq:window} is then an exact
coefficient-valuation formula for powers of the usual formal logarithm. The proof uses only
$q-1\ge1$, so it is uniform in the residue degree. The threshold has the closed finite
expression
\[
 J_p(N)=\min_{0\le b\le\floor{\log_p(N+1)}}p^b\Bigl(\floor*{\frac{N-b}{p^b}}+1\Bigr):
\]
the expression for a fixed $b$ is the first positive multiple of $p^b$ whose sum with $b$
exceeds $N$; it always satisfies the required inequality, and every minimizing $j$ is
included by taking $b=v_p(j)$. The formula \eqref{lub:eq:window} was checked by finite-ring
arithmetic modulo $p^{20}$ in $3033$ cases
\textup{(}\artifact{grh_transfer_advances_2026-11-26/checks/local/log_window_verification.json}\textup{)}.
\end{remark}

We now specialize to the setting of \cite[Section~5]{BergerSprang2025Integer}. Let
$K/\Qp$ be the unramified quadratic extension, so $q=p^2$ and $\pi=p$, with the coordinate
$[p](X)=pX+X^q$. Let $\mathcal M=\OO[[X]]^{\psi=0}$ be the kernel of the normalized trace.
By \cite[Lemma~5.2]{BergerSprang2025Integer}, every element of $\mathcal M$ is a convergent
$\OO$-linear combination of the series
\begin{equation}\label{lub:eq:generators}
 X^i(pX+X^q)^j\quad(1\le i\le q-2,\ j\ge0),\qquad
 (pX^{q-1}+q-1)(pX+X^q)^j\quad(j\ge0),
\end{equation}
which belong to $\mathcal M$. For $f\in\mathcal M$ let $c_{f,m}(T)$ be the polynomial
characterized by $c_{f,m}(x)=[X^m]f([x](X))$ for $x\in\OO$, so that
$f([T](X))=\sum_{m\ge0}c_{f,m}(T)X^m$. Put
\[
 w_q(n)=\sum_{j\ge1}\floor{n/q^j}=\frac{n-s_q(n)}{q-1},
\]
where $s_q(n)$ is the base-$q$ digit sum.

\begin{definition}[{\cite[Definition~5.1]{BergerSprang2025Integer}}]\label{lub:def:s0}
For $n\ge0$, $s_0(n)$ is the least cutoff $N\ge n$ at which the $\OO$-span of the
polynomials $c_{f,m}$ with $f\in\mathcal M$ and $m\le N$ contains a polynomial of degree
$n$ whose leading coefficient has valuation $-w_q(n)$.
\end{definition}

The finiteness of $s_0(n)$ follows from the quadratic character-variety result and the
criterion of \cite{BergerSprang2025Integer,BergerSprangSteingart2026Bounded}. Let
$E(Z)=Z+\cdots$ be the Lubin--Tate exponential, the compositional inverse of $L$, so that
$[x](X)=E(xL(X))$.

\begin{lemma}[Coefficient identity]\label{lub:lem:identity}
For $f\in\OO[[X]]$ and $m,n\ge0$,
\[
 [T^n]\,c_{f,m}(T)=[Z^n]f(E(Z))\cdot[X^m]L(X)^n .
\]
The second factor vanishes unless $m\ge n$ and $m\equiv n\bmod(q-1)$, and equals $1$ for
$m=n$. In particular $\deg c_{f,m}\le m$, and
$c_{f,m}(\zeta T)=\zeta^mc_{f,m}(T)$ for every $(q-1)$st root of unity $\zeta$. For
$f\in\mathcal M$ and $n>0$ the first factor has valuation at least $-w_q(n)$.
\end{lemma}
\begin{proof}
Substitute $[T](X)=E(TL(X))$ into $f$ and expand $f(E(Z))$ at $Z=TL(X)$. The vanishing is
the support statement of Theorem~\ref{lub:thm:window} together with $L(X)^n=X^n+\cdots$;
equivalently, $[\zeta](X)=\zeta X$ in this coordinate, because $\zeta X$ commutes with
$[p]$ and has the required derivative. Since $c_{f,n}$ is integer-valued on $\OO$ and has
degree at most $n$, its coefficient at $T^n$, which is the first factor, has valuation at
least $-w_q(n)$.
\end{proof}

\begin{corollary}[A sharp coefficient obstruction; \stU]\label{lub:cor:obstruction}
Let $n>0$. If $s_0(n)>n$, then
\[
 s_0(n)\ \ge\ n+(q-1)\,J_p\bigl(v_p(n)\bigr).
\]
In particular $s_0(n)\ge n+q-1$, and $s_0(n)\ge n+(q-1)(J+1)$ whenever $J\ge1$ satisfies
$J+\floor{\log_pJ}\le v_p(n)$.
\end{corollary}
\begin{proof}
By Lemma~\ref{lub:lem:identity}, rows $c_{f,m}$ with $m<n$ have degree $<n$, so a polynomial of
degree $n$ in the span at cutoff $n$ has leading coefficient $[Z^n]f(E(Z))$ for some
$f\in\mathcal M$. Failure at depth $n$, the bound of Lemma~\ref{lub:lem:identity} and the
integral value group of $K$ imply that $[Z^n]f(E(Z))$ has valuation at least $1-w_q(n)$
for every $f\in\mathcal M$. Theorem~\ref{lub:thm:window} shows that all second factors
$[X^m]L(X)^n$ with $m<n+(q-1)J_p(v_p(n))$ are integral. Thus all $T^n$-coefficients in the
module generated by the rows before the indicated threshold have valuation at least
$1-w_q(n)$, which excludes the required polynomial of degree $n$. For the last assertion,
every $j\le J$ has $j+v_p(j)\le J+\floor{\log_pJ}\le v_p(n)$, so $J_p(v_p(n))\ge J+1$.
\end{proof}

\begin{theorem}[First-depth criterion for small digit sum; \stU]\label{lub:thm:firstdepth}
Let $n>0$ have the base-$q$ expansion $n=\sum_{j\ge0}n_jq^j$, and put $s=\sum_jn_j$.
\begin{enumerate}[label=\textup{(\alph*)}]
\item If $1\le s\le q-2$, then $s_0(n)=n$ holds exactly when the addition of the integers
$n_j$ in base $p$ involves no carries.
\item If $s=q-1$, then $s_0(n)>n$.
\item At the first excluded digit sum, $s_0(q-1)=2q-2$ exactly.
\end{enumerate}
In the cases with $s_0(n)>n$, that is $s=q-1$, or $s\le q-2$ with a carry,
Corollary~\ref{lub:cor:obstruction} gives $s_0(n)\ge n+(q-1)J_p(v_p(n))$.
\end{theorem}
\begin{proof}
By Lemma~\ref{lub:lem:identity}, only the $n$th row can supply the leading term at cutoff $n$,
and $s_0(n)=n$ holds exactly when some $f\in\mathcal M$ has
$v\bigl([Z^n]f(E(Z))\bigr)=-w_q(n)$.

Take $\eta^{q-1}=p$ and set $V(Z)=E(\eta Z)/\eta$, the compositional inverse of the series
$\Lambda$ in the proof of Theorem~\ref{lub:thm:window}. The exponential equation becomes
$V(pZ)=pV(Z)+pV(Z)^q$. Its coefficients are integral: for $m>1$,
\[
 (p^{m-1}-1)[Z^m]V=[Z^m]V^q,
\]
with a unit on the left and only earlier coefficients on the right. In fact this recursion
puts all coefficients in $\Zp$, with support in the degrees congruent to $1$ modulo $q-1$.
Reducing modulo $p$ yields $\overline V+\overline V^{\,q}=Z$, whence
$\overline V=\sum_{j\ge0}(-1)^jZ^{q^j}$. A monomial $X^i$ of $f$ can only contribute to
$[Z^n]f(E(Z))$ if $i\equiv n\bmod(q-1)$, and since $[Z^n]E(Z)^i=\eta^{i-n}[Z^n]V(Z)^i$ its
contribution has valuation at least $-(n-i)/(q-1)$. When $1\le s\le q-1$, the first
eligible positive $i$ is $s$, and larger $i$ give valuation strictly above $-w_q(n)$. The
constant monomial contributes zero.

For $i=s$, the reduction of $\eta^{n-s}[Z^n]E^s$ is the coefficient of $Z^n$ in
$\overline V^{\,s}$. Since $s<q$, writing $n$ as a sum of $s$ powers of $q$ uses exactly
$n_j$ copies of $q^j$. The coefficient is therefore, up to sign, the multinomial
coefficient $s!/\prod_jn_j!$, which is a unit exactly when the addition of the $n_j$ is
carry-free in base $p$, by Legendre's factorial formula. If $s\le q-2$, the monomial $X^s$
belongs to $\mathcal M$ by \eqref{lub:eq:generators}; this proves (a). For $s=q-1$,
reducing the generators \eqref{lub:eq:generators} modulo $(p,X^q)$ leaves only the constant
and the powers $X^1,\dots,X^{q-2}$, so $[X^{q-1}]f\in p\OO$ for every $f\in\mathcal M$;
this proves (b).

For (c), part (b) and Corollary~\ref{lub:cor:obstruction} give $s_0(q-1)\ge2q-2$. Set $m=2q-2$.
The series $g=pX^{q-1}+q-1$ and $g'=X^{q-2}[p](X)-g=X^m+(1-q)$ both belong to
$\mathcal M$ by \eqref{lub:eq:generators}. The logarithm has coefficient $(p-p^q)^{-1}$ at
$X^q$, as follows by comparing coefficients in $L(pX+X^q)=pL(X)$. Therefore
\[
 c_{g,m}(T)=\frac{p(q-1)}{p-p^q}\,(T^{q-1}-T^m),\qquad c_{g',m}(T)=T^m .
\]
The displayed multiplier is a unit. The span of these two rows contains $T^{q-1}$, whose
leading valuation is $0=-w_q(q-1)$, at the claimed cutoff.
\end{proof}

\begin{corollary}[Unbounded excess depth; \stU]\label{lub:cor:excess}
For every prime $p$ the excess $s_0(n)-n$ is unbounded. Along $n=(q-1)q^k$ it is at least
$(q-1)J_p(2k)\ge(q-1)\bigl(2k-O(\log(k+1))\bigr)$ as $k\to\infty$.
\end{corollary}
\begin{proof}
For $n=(q-1)q^k$ the digit sum is $q-1$ and $v_p(n)=2k$, so
Theorem~\ref{lub:thm:firstdepth}(b) and Corollary~\ref{lub:cor:obstruction} apply; for $k\ge1$ take
$J=2k-\floor{\log_p(2k)}$ in the last assertion of Corollary~\ref{lub:cor:obstruction}.
\end{proof}

For $p=2$ and $q=4$, Theorem~\ref{lub:thm:firstdepth} gives $s_0(2\cdot4^j)=2\cdot4^j$ for every
$j\ge0$, and $s_0(4^i+4^j)\ge4^i+4^j+3$ for $i\ne j$, since $1+1$ carries in base two. In
particular $s_0(5)\ge8$ and $s_0(3)=6$.

The equality \eqref{lub:eq:window} determines the exact first integrality threshold of the
second factor in Lemma~\ref{lub:lem:identity}, but the bound of Corollary~\ref{lub:cor:obstruction} is
not in general the value of $s_0(n)$. At later depths a $T^n$-coefficient of the right
valuation may occur in a polynomial of higher degree, and the higher terms must then be
eliminated by integral combinations of genuine trace-kernel rows. This is a distinct
arithmetic constraint, as the following exact values show.

\begin{proposition}[Exact values of the coefficient depth; \stU]\label{lub:prop:values}
The following values of $s_0(n)$ are exact. The two valuation columns give the
leading-coefficient ideal for degree $n$, after imposing that all higher-degree terms
vanish, at the cutoffs $s_0(n)-1$ and $s_0(n)$.
\begin{center}\upshape
\begin{tabular}{rrrrrr}
\toprule
$p$ & $q$ & $n$ & $s_0(n)$ & valuation at $s_0(n)-1$ & valuation at $s_0(n)$\\
\midrule
2&4&3&6&$1$&$0$\\
2&4&5&8&$0$&$-1$\\
2&4&7&16&$0$&$-1$\\
2&4&12&30&$-2$&$-3$\\
3&9&8&16&$1$&$0$\\
\bottomrule
\end{tabular}
\end{center}
For $p=2$ and $n=12$ one has $v_2(12)=2$ and $J_2(2)=2$, so Theorem~\ref{lub:thm:window} first
allows a suitable $T^{12}$-coefficient in a row at depth $18$, and indeed the generator
$g(X)=2X^3+3$ already gives $[T^{12}]c_{g,18}(T)=15811/9815288$, of valuation
$-3=-w_4(12)$; nevertheless $s_0(12)=30$.
\end{proposition}
\begin{proof}
The values are certified by exact computations whose outputs are in the delivered files
listed at the end of the proof; the certificate stores both the obstruction immediately
before each claimed depth and the attainment at that depth. The computation starts only
from the generators \eqref{lub:eq:generators}. For a cutoff $N$ it suffices to use
$j\le N$ in the second family and $i+j\le N$ in the first: the omitted generators have
$X$-adic order greater than $N$, since $pX+X^q$ has $X$-adic order one, and the
coefficients through $X^N$ of any element of $\mathcal M$ depend only on this finite list.
The rows $c_{f,m}$, $m\le N$, are formed exactly over $\Q$ from the Carleman matrix
$\mathcal D$ of the logarithm as
$\mathcal D^{-1}\operatorname{diag}(1,T,T^2,\dots)\mathcal D$, and all of them are
retained, including diagonal rows with nonunit coefficients. This proves completeness at
each finite cutoff.

To read off the leading-coefficient ideals, clear denominators in each rational row using
only $p$-adic units and a common factor $p^{\kappa_0}$. Choose a triangular full-rank
source submatrix whose determinant has valuation $\mu$. Its adjugate proves that
$p^{\mu}\Zp^{\,\ell}$, with $\ell$ the number of columns, is already contained in the
scaled row module, so adding $p^{\mu}I$ and reducing entries modulo $p^{\mu}$ preserves
that module. An exact integer column Hermite normal form now gives its leading-degree
ideals, and subtracting $\kappa_0$ restores the original valuations. Extension of scalars
to $\OO$ preserves the calculation. Monotonicity of the generated modules in the cutoff
shows that the tabulated values are exact minima and not merely upper bounds from one
witness. An independent descending-degree elimination over $\Z_{(p)}$, in which a pivot is
replaced only when a smaller valuation is found and every row operation is invertible,
gives the same values. The files are:
\begin{itemize}[nosep]
\raggedright
\item the Hermite-normal-form certificate and its generator:
\artifact{grh_transfer_advances_2026-11-26/checks/local/depth_hnf_checks.json},
\artifact{grh_transfer_advances_2026-11-26/checks/local/verify_depth_hnf.py};
\item the elimination output for $p=2$, cutoff $48$, and for $p=3$, cutoff $100$, where a
null entry means only that the cutoff was exceeded:
\artifact{grh_transfer_advances_2026-11-26/checks/local/depth_p2_48.json},
\artifact{grh_transfer_advances_2026-11-26/checks/local/depth_p3_100.json}.\qedhere
\end{itemize}
\end{proof}

\begin{remark}[The numerical baseline]\label{lub:rem:baseline}
The values in Proposition~\ref{lub:prop:values} refer to Definition~\ref{lub:def:s0} as stated. According to
Berger's errata \cite[Section~19]{Berger2026Errata}, the implementation in Appendix~A of
\cite{BergerSprang2025Integer} inserts $T^s$ at each stage without establishing membership
in the module and omits the possible nonunit diagonal generator, so that the printed
experiment there addresses $\operatorname{Pol}(\mathcal M)+\OO[T]$ rather than the defined
$s_0(n)$. The definition, the generators \eqref{lub:eq:generators} and the proofs above are
not affected.
\end{remark}

\subsection{Open questions}\label{lub:ssec:open}
\begingroup\raggedright\sources{grh_transfers_cycle_signings_and_lubin_tate; grh_transfers_lubin_tate_and_graph_towers; grh_prime_patterns_and_branch_precision; grh_method_transfers_orthogonal_jets_2026-11-26; grh_transfer_advances_2026-11-26}\endgroup

\begin{question}\label{lub:q:closed}
When the carry criterion \eqref{lub:eq:carry} fails and $h_*\ge2$, the defect $\eta_r$ of
Theorem~\ref{lub:thm:defect} is computed by a finite recurrence and depends on the shallower
contacts. Is there a closed formula for $\eta_r$ in terms of the contact multiplicities
$\nu_j$, in particular over ramified fields?
\end{question}

\begin{question}\label{lub:q:period}
Determine the least eventual period of $r\mapsto t_b(r)-Dr$ directly from the contact
multiplicities. The periods $\Theta$ of Theorem~\ref{lub:thm:defect} and $p^{\vartheta}$ of
Corollary~\ref{lub:cor:depthoneperiod} are explicit but not claimed to be minimal.
\end{question}

\begin{question}\label{lub:q:growth}
What is the true growth of $s_0\bigl((q-1)q^k\bigr)-(q-1)q^k$? Corollary~\ref{lub:cor:excess}
supplies a growing lower bound. A matching upper bound would require the explicit
construction of trace-kernel rows at a controlled depth, together with the integral
elimination of the higher-degree terms; by Proposition~\ref{lub:prop:values} the coefficient threshold
of Theorem~\ref{lub:thm:window} does not by itself control that cancellation.
\end{question}

\begin{question}\label{lub:q:delayed}
Determine $s_0(n)$ exactly for the delayed families of Theorem~\ref{lub:thm:firstdepth}, and
describe its behaviour for digit sums $s\ge q$, which the first-depth criterion does not
cover.
\end{question}
